\documentclass[12pt,a4paper]{article}
\usepackage[T1]{fontenc}
\usepackage{lmodern,amsmath,amssymb}
\usepackage[margin=22mm,top=18mm,bottom=19mm]{geometry}
\usepackage{microtype}
\pdfinfoomitdate=1
\pdftrailerid{}
\pdfsuppressptexinfo=15
\usepackage[hidelinks,pdflang={en},
  pdfauthor={Astra 6.2},
  pdftitle={A nonfree factor of the binary free monoid},
  pdfsubject={Supervised by Christophe Cordero. Discovered on 25 November 2026.}]{hyperref}
\pagestyle{empty}
\setlength{\parindent}{0pt}
\setlength{\parskip}{.65em}
\newcommand{\w}[1]{\mathit{#1}}
\begin{document}
{\LARGE A nonfree factor of the binary free monoid\par}
\smallskip
\begin{tabular*}{\linewidth}{@{}l@{\extracolsep{\fill}}ll@{}}
{\small\sffamily Author} & {\small\sffamily Supervisor} & {\small\sffamily Date of discovery} \\
Astra 6.2 & Christophe Cordero & 25 November 2026
\end{tabular*}
\par\smallskip\hrule\medskip

\textbf{Theorem.} Let $A^*$ be the monoid of finite words on $A=\{a,b\}$
under concatenation, and let $t^*$ denote all powers of a word $t$, including the
empty word. There is a nonfree submonoid $M$ of $A^*$ such that every word has
a unique decomposition in the product
\[
  A^*=(\w{ab})^*\,b^*\,M\,a^*\,(\w{aabb})^*.
\]

\textit{Proof.} Let $D$ be the words without prefix $\w{ab}$, and $E$ those
without suffix $\w{aabb}$. Let $L$ consist of the words in $D$ whose initial
$b$, if present, is followed by $\w{ab}$. Let $R$ consist of the words in $E$
whose final $a$, if present, is preceded by $\w{aabb}$. Set $M=L\cap R$.

Taking the longest initial power of $\w{ab}$ gives $A^*=(\w{ab})^*D$ uniquely.
Taking the longest terminal power of $\w{aabb}$ likewise gives
$A^*=E(\w{aabb})^*$. The definitions give the disjoint unions
$D=L\sqcup bD$ and $E=R\sqcup Ea$. In the second alternative of either union,
deleting the indicated letter leaves a shorter word in the same language.
Induction on length therefore gives the unambiguous products
\[
  A^*=(\w{ab})^*\,b^*\,L, \qquad A^*=R\,a^*\,(\w{aabb})^*.
\]
If appending $s\in a^*(\w{aabb})^*$ to a word $x$ created a prefix $b$,
$\w{ab}$, or $\w{bab}$, then $s$ would start with a nonempty suffix of one of
these three words. The only such suffixes are $b$, $\w{ab}$, and $\w{bab}$.
None can begin $s$, which is empty, equals $a$, or starts with $\w{aa}$.
Thus $xs\in L$ exactly when $x\in L$. The second decomposition restricts to
$L=M\,a^*\,(\w{aabb})^*$. Substitution in the first gives the required unique
decomposition.

The empty word belongs to $M$. The only words of length one or two in $L$ are
$a$ and $\w{aa}$, and both are excluded by $R$. Thus nonempty $u,v\in M$ have
length at least three. The first three letters of $uv$ come from $u$, so
$uv\in L$. If $v$ ends in $a$, it already ends in $\w{aabba}$. A final
$\w{aabb}$ in $uv$ would have to cross the join, forcing $v=\w{abb}$, which
violates the prefix rule. Hence $uv\in M$.

The four parenthesized words in
\[
  (\w{aab})(\w{babab})=\w{aabbabab}=(\w{aabba})(\w{bab})
\]
lie in $M$ because each starts with $\w{aa}$ or $\w{bab}$ and ends with
$\w{ab}$ or $\w{aabba}$. Each has length three or five, so none splits into two
nonempty elements of $M$. If $M$ were free, each would be a letter of its free
generating alphabet, since any longer word over that alphabet splits after its
first letter. Equality would then force $\w{aab}=\w{aabba}$, contradicting
their lengths.\hfill$\square$
\end{document}
